% !TeX program = lualatex
\RequirePackage{iftex}
\ifLuaTeX\else
  \GenericError{}{This document must be compiled with LuaLaTeX.}{}{Run `latexmk -lualatex sample.tex` or set your editor's TeX engine to LuaLaTeX.}
  \expandafter\endinput
\fi
\documentclass[9pt]{beamer}
\usepackage{luatexja}
\usepackage{mathtools}
\usepackage{amssymb}
\usepackage{amsthm}
\usepackage{braket}
\usepackage{bm}
\usepackage{booktabs}
\usepackage{array}

\setbeamertemplate{proof begin}{}
\setbeamertemplate{proof end}{}

\usetheme[block=fill]{metropolis}

%%%%%%%%%%%%%%%%%%%%%%%%%%%%%%%% １ページ目 %%%%%%%%%%%%%%%%%%%%%%%%%%%%%%%%%
\title{数学講究第２回}
\author{今泉 力}
\date{2026年4月29日}

\begin{document}
\begin{frame}
  \titlepage
\end{frame}

%%%%%%%%%%%%%%%%%%%%%%%%%%%%%%%% ２ページ目 %%%%%%%%%%%%%%%%%%%%%%%%%%%%%%%%%
\begin{frame}
  \frametitle{Sample Frame}
  This is a sample frame in a Beamer presentation.
  \begin{enumerate}
    \item First item
    \item Second item
    \item Third item
  \end{enumerate}
  \begin{itemize}
    \item Bullet point one
    \item Bullet point two
    \item Bullet point three
  \end{itemize}
\end{frame}

%%%%%%%%%%%%%%%%%%%%%%%%%%%%%%%% ３ページ目 %%%%%%%%%%%%%%%%%%%%%%%%%%%%%%%%%
\begin{frame}
  \frametitle{block環境のサンプル}
  \begin{block}{Block Title}
    これはblock環境のサンプルです。
  \end{block}
  \begin{alertblock}{Alert Block Title}
    これはalertblock環境のサンプルです。
  \end{alertblock}
  \begin{exampleblock}{Example Block Title}
    これはexampleblock環境のサンプルです。
  \end{exampleblock}
\end{frame}

%%%%%%%%%%%%%%%%%%%%%%%%%%%%%%%% ４ページ目 %%%%%%%%%%%%%%%%%%%%%%%%%%%%%%%%%

\begin{frame}
  \frametitle{n量子ビットの定義}
  \begin{block}{n量子ビットの定義}
    以下の集合$Q_{n}$の要素を\textbf{n量子ビット}と呼ぶ。
    \[
    Q_n \coloneq \left\{ \begin{pmatrix}a_1 \\ a_2 \\ \vdots \\ a_{2^n} \end{pmatrix} \in \mathbb{C}^{2^n} \middle| \sum_{i=1}^{2^n} |a_i|^2 = 1 \right\}
    \]
  \end{block}
  \begin{block}{ブラケット表記}
    \begin{itemize}
      \item 量子力学において、$Q_n$の要素は、$\bm{\ket{\varphi}}$のように表されることが多い。これを\textbf{ケット・ベクトル}と呼ぶ。
      \item また、ある$\bm{\ket{\psi}} \in Q_n$に対して、その共役転置（随伴）を$\bm{\bra{\psi}}$と表し、これを\textbf{ブラ・ベクトル}と呼ぶ。
    \end{itemize}
  \end{block}
\end{frame}

%%%%%%%%%%%%%%%%%%%%%%%%%%%%%%%%% ５ページ目 %%%%%%%%%%%%%%%%%%%%%%%%%%%%%%%%%
\begin{frame}
  \frametitle{量子ビット（$\ket{0}$と$\ket{1}$の定義）}
  \begin{block}{量子ビット（$\ket{0}$と$\ket{1}$の定義）}
    \begin{itemize}
      \item 量子ビット$\ket{\varphi} \in Q_1$は、以下のように表される。
      \[\ket{\varphi} = \alpha \ket{0} + \beta \ket{1} \quad (\alpha, \beta \in \mathbb{C}, |\alpha|^2 + |\beta|^2 = 1)\]
      ただし、$\ket{0} \coloneq \begin{pmatrix}1 \\ 0\end{pmatrix}$、$\ket{1} \coloneq \begin{pmatrix}0 \\ 1\end{pmatrix}$と定義する。
    \end{itemize}
  \end{block}
\end{frame}

%%%%%%%%%%%%%%%%%%%%%%%%%%%%%%%% ６ページ目 %%%%%%%%%%%%%%%%%%%%%%%%%%%%%%%%%
\begin{frame}
  \frametitle{内積の定義}
  \begin{block}{随伴の定義}
    $\bm{a} \in \mathbb{C}^{n}$ に対して、その随伴（共役転置）を $\bm{a}^\dagger$ と表す。
  \end{block}
  \begin{block}{内積の定義}
    $\bm{a}, \bm{b} \in \mathbb{C}^{n}$ に対して、内積は以下のように定義される。
    \[(\bm{a}, \bm{b}) \coloneq \bm{a}^\dagger \bm{b} = \sum_{i=1}^{n} a_i^* b_i\]
    ブラケット表記では以下のように書く。（$\bm{\ket{\varphi}}, \bm{\ket{\psi}} \in Q_n$とする。）
    \[\braket{\varphi | \psi} \coloneq \bm{\ket{\varphi}}^\dagger \bm{\ket{\psi}}\]
  \end{block}
\end{frame}

%%%%%%%%%%%%%%%%%%%%%%%%%%%%%%%% ７ページ目 %%%%%%%%%%%%%%%%%%%%%%%%%%%%%%%%%
\begin{frame}
  \frametitle{ユニタリ行列}
  \begin{block}{ユニタリ行列の定義}
    $U \in Mat_{n}(\mathbb{C})$がユニタリ行列であるとは、以下の条件を満たすことをいう。
    \[U^\dagger U = U U^\dagger = I_n\]
    ここで、$I_n$は$n$次の単位行列である。
  \end{block}
  \begin{block}{定理：ユニタリ行列と量子ビットの関係}
    $$
    \forall \ket{\varphi} \in Q_n: Uはユニタリ行列 \Rightarrow U\ket{\varphi} \in Q_n
    $$
  \end{block}
  \begin{proof}
    \textbf{proof.} \\
    $(U\ket{\varphi})^\dagger (U\ket{\varphi}) = \bra{\varphi} (U^\dagger U) \ket{\varphi} = \braket{\varphi | \varphi} = 1$
  \end{proof}
\end{frame}

%%%%%%%%%%%%%%%%%%%%%%%%%%%%%%%% ８ページ目 %%%%%%%%%%%%%%%%%%%%%%%%%%%%%%%%%
\begin{frame}
  \frametitle{テンソル積の定義}
  \begin{block}{テンソル積の定義}
    \begin{itemize}
      \item $\begin{pmatrix} a_1 \\ \vdots \\ a_{2^n} \end{pmatrix}, \begin{pmatrix}b_1 \\ \vdots \\ b_{2^n} \end{pmatrix} \in Q_n$に対して、テンソル積は以下のように定義される。
      \[\ket{a}\otimes \ket{b} = \begin{pmatrix}a_1\\ \vdots \\ a_{2^n} \end{pmatrix} \otimes \begin{pmatrix}b_1\\ \vdots \\ b_{2^n} \end{pmatrix} \coloneq \begin{pmatrix}a_1 \begin{pmatrix}b_1 \\ \vdots \\ b_{2^n} \end{pmatrix} \\ \vdots \\ a_{2^n} \begin{pmatrix}b_1 \\ \vdots \\ b_{2^n} \end{pmatrix} \end{pmatrix} = \begin{pmatrix}a_1 b_1 \\ a_1 b_2 \\ \vdots \\ a_1 b_{2^n} \\ a_2 b_1 \\ \vdots \\ a_{2^n} b_{2^n} \end{pmatrix}\]
      
      \item $\ket{a}\otimes \ket{b}$は$\ket{ab}$とも書かれる。
    \end{itemize}
  \end{block}
\end{frame}

%%%%%%%%%%%%%%%%%%%%%%%%%%%%%%%% ９ページ目 %%%%%%%%%%%%%%%%%%%%%%%%%%%%%%%%%
\begin{frame}
  \frametitle{テンソル積の性質}
  \begin{block}{定理：テンソル積の性質}
    $\forall \ket{\varphi}, \ket{\psi}, \ket{\chi}, \ket{\omega} \in Q_n, \forall c \in \mathbb{C}$に対して、以下が成り立つ。
    \begin{enumerate}
      \item $(c\ket{\varphi}) \otimes \ket{\psi} = c(\ket{\varphi} \otimes \ket{\psi} )$
      \item $\ket{\varphi} \otimes (c\ket{\psi}) = c(\ket{\varphi} \otimes \ket{\psi} )$
      \item $(\ket{\varphi} + \ket{\psi}) \otimes \ket{\chi} = \ket{\varphi} \otimes \ket{\chi} + \ket{\psi} \otimes \ket{\chi}$
      \item $\left( \bra{\varphi} \otimes \bra{\psi} \right) \left( \ket{\chi} \otimes \ket{\omega} \right) = \braket{\varphi | \chi} \braket{\psi | \omega}$
    \end{enumerate}
  \end{block}
  \begin{proof}
    \textbf{proof.} \\
    1,2$\dots$スカラーの積は可換より \\
    3$\dots$スカラーは分配法則が成り立つため。
    4$\dots$左辺と右辺を別々で計算すると成り立つ。
  \end{proof}
\end{frame}

%%%%%%%%%%%%%%%%%%%%%%%%%%%%%%%%% １０ページ目 %%%%%%%%%%%%%%%%%%%%%%%%%%%%%%%%%
\begin{frame}
  \frametitle{量子ゲート（代表的な１量子ビットゲート）}
  \begin{block}{量子ゲート}
      量子コンピュータにおいて、n量子ビットに対して作用するユニタリ行列を\textbf{n量子ビットゲート}と呼ぶ。
  \end{block}
  \begin{table}[]
    \centering
    \footnotesize
    \begin{tabular}{|c|>{\centering\arraybackslash}p{5em}|c|c|} \hline
      ゲート & 意味 & 量子回路 & 行列 \\ \hline
      Xゲート & NOTゲート & \hspace{5em} \rule{0pt}{2em} & $\begin{bmatrix}0 & 1 \\ 1 & 0\end{bmatrix}$ \\ \hline
      Yゲート & 位相反転ゲート & \hspace{5em} \rule{0pt}{2em} & $\begin{bmatrix}0 & -i \\ i & 0\end{bmatrix}$ \\ \hline
      Zゲート & 位相反転ゲート & \hspace{5em} \rule{0pt}{2em} & $\begin{bmatrix}1 & 0 \\ 0 & -1\end{bmatrix}$ \\ \hline
      Hゲート & アダマールゲート & \hspace{5em} \rule{0pt}{2em} & $\frac{1}{\sqrt{2}} \begin{bmatrix}1 & 1 \\ 1 & -1\end{bmatrix}$ \\ \hline
      Sゲート & 位相ゲート & \hspace{5em} \rule{0pt}{2em} & $\begin{bmatrix}1 & 0 \\ 0 & i\end{bmatrix}$ \\ \hline
      Tゲート & 位相ゲート & \hspace{5em} \rule{0pt}{2em} & $\begin{bmatrix}1 & 0 \\ 0 & e^{i\pi/4}\end{bmatrix}$ \\ \hline
    \end{tabular}
  \end{table}
\end{frame}

%%%%%%%%%%%%%%%%%%%%%%%%%%%%%%%%% １１ページ目 %%%%%%%%%%%%%%%%%%%%%%%%%%%%%%%%%
\begin{frame}
  \frametitle{量子ゲート（特別な記法のある量子ゲート）}
  \begin{table}[]
    \centering
    \scriptsize
    \begin{tabular}{|c|>{\arraybackslash}p{12em}|c|c|} \hline
      ゲート & 意味 & 量子回路 & 行列 \\ \hline
      CNOT & コントロール量子ビットが$\ket{1}$のとき、ターゲット量子ビットを反転する & \hspace{7em} \rule{0pt}{1em} & $\begin{bmatrix}1 & 0 & 0 & 0 \\ 0 & 1 & 0 & 0 \\ 0 & 0 & 0 & 1 \\ 0 & 0 & 1 & 0\end{bmatrix}$ \\ \hline
      SWAP & ２つの量子ビットの交換 & \hspace{7em} \rule{0pt}{1em} & $\begin{bmatrix}1 & 0 & 0 & 0 \\ 0 & 0 & 1 & 0 \\ 0 & 1 & 0 & 0 \\ 0 & 0 & 0 & 1\end{bmatrix}$ \\ \hline
      制御Z & コントロール量子ビットが$\ket{1}$のとき、ターゲット量子ビットを位相反転する & \hspace{7em} \rule{0pt}{1em} & $\begin{bmatrix}1 & 0 & 0 & 0 \\ 0 & 1 & 0 & 0 \\ 0 & 0 & 1 & 0 \\ 0 & 0 & 0 & -1\end{bmatrix}$ \\ \hline
      Toffoli & トリプルコントロールNOTゲート & \hspace{7em} \rule{0pt}{1em} & $\begin{bmatrix}1 & 0 & 0 & 0 & 0 & 0 & 0 & 0 \\ 0 & 1 & 0 & 0 & 0 & 0 & 0 & 0 \\ 0 & 0 & 1 & 0 & 0 & 0 & 0 & 0 \\ 0 & 0 & 0 & 1 & 0 & 0 & 0 & 0 \\ 0 & 0 & 0 & 0 & 1 & 0 & 0 & 0 \\ 0 & 0 & 0 & 0 & 0 & 1 & 0 & 0 \\ 0 & 0 & 0 & 0 & 0 & 0 & 0 & 1 \\ 0 & 0 & 0 & 0 & 0 & 0 & 1 & 0\end{bmatrix}$ \\ \hline
    \end{tabular}
  \end{table}
\end{frame}

%%%%%%%%%%%%%%%%%%%%%%%%%%%%%%%%%% １２ページ目 %%%%%%%%%%%%%%%%%%%%%%%%%%%%%%%%%
\begin{frame}
  \frametitle{量子回路の書き方}
\end{frame}

%%%%%%%%%%%%%%%%%%%%%%%%%%%%%%%%%%% １３ページ目 %%%%%%%%%%%%%%%%%%%%%%%%%%%%%%%%%
\begin{frame}
	\frametitle{量子回路の書き方}
    まず、初期状態に$\ket{000}$にアダマールゲートを作用させる。
  \begin{align}
    \ket{000} &\xrightarrow{H_1, H_3} \left(\frac{1}{\sqrt{2}}(\ket{0} + \ket{1}) \right) \ket{0} \left(\frac{1}{\sqrt{2}}(\ket{0} + \ket{1}) \right) \\
    &= \frac{1}{2}(\ket{000} + \ket{001} + \ket{100} + \ket{101})
  \end{align}
  次に、第１量子ビットをコントロール量子ビットとするCNOTゲートを第２量子ビットに作用させる。
  \begin{align}
    \xrightarrow{CNOT_{1,2}} \frac{1}{2}(\ket{000} + \ket{001} + \ket{110} + \ket{111})
  \end{align}
  第２量子ビットをコントロール量子ビットとするCZゲートを第３量子ビットに作用。Z$\ket{0} = \ket{0}$、Z$\ket{1} = -\ket{1}$より、
  \begin{align}
    \xrightarrow{CZ_{2,3}} \frac{1}{2}(\ket{000} + \ket{001} + \ket{110} - \ket{111})
  \end{align}
\end{frame}

%%%%%%%%%%%%%%%%%%%%%%%%%%%%%%%%%% １４ページ目 %%%%%%%%%%%%%%%%%%%%%%%%%%%%%%%%%
\begin{frame}
  \frametitle{量子回路の書き方}
  最後に、アダマールゲートを第３量子ビットに作用させる。H$\ket{0} = \frac{1}{\sqrt{2}}(\ket{0} + \ket{1})$、H$\ket{1} = \frac{1}{\sqrt{2}}(\ket{0} - \ket{1})$より、
  \begin{align}
    &\xrightarrow{H_3} \frac{1}{2} \Bigl( \ket{00}\frac{1}{\sqrt{2}}(\ket{0} + \ket{1}) + \ket{00}\frac{1}{\sqrt{2}}(\ket{0} - \ket{1}) \\
    &\qquad + \ket{11}\frac{1}{\sqrt{2}}(\ket{0} + \ket{1}) - \ket{11}\frac{1}{\sqrt{2}}(\ket{0} - \ket{1}) \Bigr) \\
    &= \frac{1}{2\sqrt{2}} \left( \ket{000} +\ket{001} + \ket{000} - \ket{001} + \ket{110} + \ket{111} - \ket{110} + \ket{111} \right) \\
    &= \frac{1}{\sqrt{2}}(\ket{000} + \ket{111})
  \end{align}
\end{frame}

\begin{frame}
  \frametitle{量子もつれ（エンタングルメント）}
  \begin{block}{量子もつれ（エンタングルメント）}
    量子ビット$\ket{\varphi} \in Q_n$が、ある$\ket{\psi} \in Q_{n-1}$と$\ket{\chi} \in Q_1$を用いて$\ket{\varphi} = \ket{\psi} \otimes \ket{\chi}$と表せないとき、$\ket{\varphi}$は\textbf{量子もつれ（エンタングルメント）状態}であるという。
  \end{block}
  例 \dots ベル状態$\ket{B_{00}} \coloneq \frac{1}{\sqrt{2}}(\ket{00} + \ket{11})$は量子もつれ状態である。
  \begin{proof}
    \textbf{proof.} \\
    $\ket{B_{00}}$が量子もつれ状態でないと仮定する。すなわち、$\ket{\psi} = a\ket{0} + b\ket{1} \in Q_1$、$\ket{\chi} = c\ket{0} + d\ket{1} \in Q_1$が存在して、$\ket{B_{00}} = \ket{\psi} \otimes \ket{\chi}$であるとする。（$a, b, c, d \in \mathbb{C}$で、$|a|^2 + |b|^2 = 1$、$|c|^2 + |d|^2 = 1$を満たす。）このとき、
    \begin{align}
      \frac{1}{\sqrt{2}}(\ket{00} + \ket{11}) &= (a\ket{0} + b\ket{1}) \otimes (c\ket{0} + d\ket{1}) \\
      &= ac\ket{00} + ad\ket{01} + bc\ket{10} + bd\ket{11}
    \end{align}
    ここで、$ad = 0$、$bc = 0$、$ac = bd$を満たす必要がある。これらの条件を満たす$a, b, c, d$は存在しないため、$\ket{B_{00}}$は量子もつれ状態である。
  \end{proof}
\end{frame}

%%%%%%%%%%%%%%%%%%%%%%%%%%%%%%%%%%% １５ページ目 %%%%%%%%%%%%%%%%%%%%%%%%%%%%%%%%%
\begin{frame}
  \frametitle{量子もつれ（エンタングルメント）}
  量子ビットAとBがベル状態$\ket{B_{00}}$にある量子ビットの間には、次のような相関がある。
  \begin{itemize}
    \item 量子ビットAを測定した結果0が得られたとき、量子ビットBを測定した結果も0になる。
    \item 量子ビットAを測定した結果1が得られたとき、量子ビットBを測定した結果も1になる。
    \item 量子ビットBを測定した結果0が得られたとき、量子ビットAを測定した結果も0になる。
    \item 量子ビットBを測定した結果1が得られたとき、量子ビットAを測定した結果も1になる。
  \end{itemize}
  ベル状態は全部で、
\begin{align}
  \ket{B_{00}} &\coloneq \frac{1}{\sqrt{2}}(\ket{00} + \ket{11}), \ket{B_{01}} \coloneq \frac{1}{\sqrt{2}}(\ket{00} - \ket{11})\\
   \ket{B_{10}} &\coloneq \frac{1}{\sqrt{2}}(\ket{01} + \ket{10}), \ket{B_{11}} \coloneq \frac{1}{\sqrt{2}}(\ket{01} - \ket{10})
\end{align}
の４つがある。これらはすべて量子もつれ状態である。
\end{frame}

%%%%%%%%%%%%%%%%%%%%%%%%%%%%%%%%%%%% １６ページ目 %%%%%%%%%%%%%%%%%%%%%%%%%%%%%%%%%
\begin{frame}
  \frametitle{量子複製不可能定理}
  \begin{block}{量子複製不可能定理}
    任意の量子ビット$\ket{\varphi} \in Q_n$に対して、$\ket{\varphi}$を複製する、つまり
    \[U\ket{\varphi}\ket{0} = \ket{\varphi}\ket{\varphi}\]
    となるユニタリ行列Uは存在しない。
  \end{block}
  \begin{proof} 
    \textbf{proof.}
    $\forall \ket{\varphi}, \ket{\psi} \in Q_n$に対して、
    \begin{equation}
      \begin{cases}
        U\ket{\varphi}\ket{0} = \ket{\varphi}\ket{\varphi} \\
        U\ket{\psi}\ket{0} = \ket{\psi}\ket{\psi}
      \end{cases}
    \end{equation}
    が成り立つと仮定する。2つの式の右辺の内積をとると、
    $$\left( \bra{\varphi} \bra{0} \right) U^\dagger U \left( \ket{\psi} \ket{0} \right) = \braket{\varphi | \psi} \quad (\because U^\dagger U = I, \text{テンソル積の性質4})$$ 
    一方、左辺の内積をとると、
    $$\left( \bra{\varphi} \bra{\varphi} \right) \left( \ket{\psi} \ket{\psi} \right) = \braket{\varphi | \psi}^2$$
    $$\therefore \braket{\varphi | \psi} = \braket{\varphi | \psi}^2$$ 
    $$\iff \braket{\varphi | \psi} = 0 \text{ or } 1$$
    よって、任意の量子状態を複製することはできない。
  \end{proof}
\end{frame}

\begin{frame}
  \frametitle{量子コンピュータにデータを入力する方法}
  \古典的なデータを量子コンピュータで処理するためには、古典データを「量子状態」に変換して入力する必要がある。これを\textbf{エンコーディング}と呼ぶ。\\
  島田義皓著「量子コンピューティング」では、以下のようなエンコーディング方法が紹介されている。
\begin{itemize}
  \item \textbf{ディジタル入力（基底エンコーディング）}：古典データのビット列をそのまま量子ビットの計算基底$\ket{0}$と$\ket{1}$に対応させる方法。\\
  \qquad 例：古典データ「5」を量子状態$\ket{101}$にエンコードする。
  \item \textbf{アナログ入力（振幅エンコーディング）}：
  \item \textbf{ハミルトニアンエンコーディング}：
\end{frame}

\end{document}  
